% Brief (grammar: every state as a point, every move as a stroke; atlas
%   VI/fig:bonded-moves).
% Reader question: from where a project's funds are, where can a move take
% them?
% Claim: no move the kernel makes leaves the triangle of funds that add up
% to what was topped up. A post or a refund moves along a row, a slash moves
% up a row, and nothing moves down. With the check of the wallet removed,
% moves leave the triangle.
% Evidence: every move made from a state with all 6 units topped up, in the
% walk of every state within the bounds [internal, harbor states --each;
% sessions states-each.txt and states-each-overdraw.txt, drawn as printed,
% and tests/harbor-research/test_bonded_ledger_numbers.py holds the counts
% in the caption to them].
% Must distinguish: a move along a row from a move up; a move the Book's
% rules allow from one they refuse; inside the triangle from outside it.
% Grammar chosen: escrow across and the pool up, so that the wallet is what
% is left and the funds that add up are a triangle. One stroke for each
% kind of move the walk made.
% Grammar rejected: an equilateral plot of three shares, where a row of
% equal pool is no longer level and the moves out of the triangle have
% nowhere to go; the table itself, which is 211 rows.
% Five-second test: on which side does a stroke cross the dotted line?
\begin{figure}[H]
\centering
\pdfigurehue{pdcobalt}%
\begin{tikzpicture}[pd figure,x=.5cm,y=.5cm,
    row move/.style={draw=pdinkmuted!75,line width=.5pt},
    rising move/.style={draw=pdfocus,line width=.35pt,-{Stealth[length=2.2pt,width=1.8pt]},shorten >=1.6pt},
    stray move/.style={draw=pderror,line width=.7pt,-{Stealth[length=3pt,width=2.6pt]},shorten >=1.8pt},
    place/.style={circle,fill=pdink,inner sep=0pt,minimum size=2.6pt},
    beyond/.style={circle,draw=pderror,fill=pdpage,line width=.6pt,inner sep=0pt,minimum size=3.2pt}]
  % The tables are the recorded runs, read as they were printed: the first
  % line is the command, and every remark starts with '#'.
  \pgfplotstableread[skip first n=1]{excerpts/states-each.txt}\pdmovesreal
  \pgfplotstableread[skip first n=1]{excerpts/states-each-overdraw.txt}\pdmovesloose
  % Counts taken from the tables as they are drawn.
  \ifdefined\pdmovesin\else \newcount\pdmovesin \newcount\pdmovesout \fi
  % One panel: the moves of a table, then the places on top of them.
  \newcommand{\pdmovespanel}[1]{%
    \global\pdmovesin=0 \global\pdmovesout=0
    \pgfplotstablegetrowsof{#1}
    \pgfmathtruncatemacro{\pdlastmove}{\pgfplotsretval-1}
    % Moves along a row first, so that the moves up lie over them.
    \foreach \pass in {0,1,2}{
      \foreach \i in {0,...,\pdlastmove}{
        \pgfplotstablegetelem{\i}{e0}\of#1 \let\ea\pgfplotsretval
        \pgfplotstablegetelem{\i}{p0}\of#1 \let\pa\pgfplotsretval
        \pgfplotstablegetelem{\i}{e1}\of#1 \let\eb\pgfplotsretval
        \pgfplotstablegetelem{\i}{p1}\of#1 \let\pb\pgfplotsretval
        \pgfplotstablegetelem{\i}{wrong}\of#1 \let\wrong\pgfplotsretval
        \ifnum\wrong=1
          \ifnum\pass=2 \global\advance\pdmovesout 1
            \draw[stray move] (\ea,\pa) -- (\eb,\pb); \node[beyond] at (\eb,\pb) {};\fi
        \else\ifnum\pa=\pb
          \ifnum\pass=0 \global\advance\pdmovesin 1 \draw[row move] (\ea,\pa) -- (\eb,\pb);\fi
        \else
          \ifnum\pass=1 \global\advance\pdmovesin 1 \draw[rising move] (\ea,\pa) -- (\eb,\pb);\fi
        \fi\fi
      }
    }
    % Every place where the funds add up to 6.
    \foreach \p in {0,...,6}{
      \pgfmathtruncatemacro{\most}{6-\p}
      \foreach \e in {0,...,\most}{ \node[place] at (\e,\p) {}; }
    }
    % Past this line a project would hold more than was topped up.
    \draw[pd guide] (6.9,-.4) -- (-.4,6.9);
    \foreach \n in {0,2,4,6}{
      \node[pd label,text=pdinkmuted,anchor=north] at (\n,-.55) {\n};
      \node[pd label,text=pdinkmuted,anchor=east] at (-.55,\n) {\n};
    }
    \node[pd label,anchor=north] at (3,-1.5) {units in escrow};
  }
  \pdmovespanel{\pdmovesreal}
  \node[pd label,rotate=90,anchor=south] at (-1.6,3) {units in the pool};
  % Each panel's words sit above it, clear of the strokes and of the
  % other panel.
  \node[pd title,anchor=south west] at (-.6,8.9) {the kernel's keeper};
  \node[pd label,anchor=north west,align=left] at (-.6,8.9)
    {\the\pdmovesin\ kinds of move,\\none past the line};
  \begin{scope}[xshift=4.4cm]
    \pdmovespanel{\pdmovesloose}
    \node[pd title,anchor=south west] at (-.6,8.9) {with the check of the wallet removed};
    \node[pd label,anchor=north west,align=left] at (-.6,8.9)
      {\the\pdmovesout\ more, each to a place\\with more than 6 units};
  \end{scope}
\end{tikzpicture}
\caption{\textbf{No move the kernel makes takes a project's funds past what
was topped up.} A project has 6 units, and the wallet holds whatever is not
in escrow or in the pool. Grey strokes are bonds posted and refunded, which
move along a row. Arrows are slashes, which move up a row and never down.
On the right the core is built without the guard of
Listing~\ref{alg:bonded-build-covered}, and each red arrow posts a bond the
wallet cannot cover. [internal, \texttt{harbor states -{}-each}]}
\label{fig:bonded-moves}
\end{figure}
